% Auto-generated by figs.py.  Do not edit; regenerate.
% Each macro is the set of numbers its figure computed while being drawn.
\providecommand{\fignote}[1]{\csname fignote@#1\endcsname}

\expandafter\def\csname fignote@A1\endcsname{The vertical axis is exaggerated 4$\times$ relative to the horizontal, so that the 1 nm interfacial oxide is visible; both scale bars are drawn. Total gate stack 31 nm. N$_{sub}$ = 1e+16 cm$^{-3}$, N$_{sd}$ = 5e+19 cm$^{-3}$, gate work function 4.35 eV. Dashed lines are the metallurgical junctions, which sit 15 nm inside each gate edge.}
\expandafter\def\csname fignote@A5\endcsname{every run executed at A = 0.071; the correction to 0.045 is a pure post-multiplier on contact quantities, an exact rescale of 0.6338$\times$ verified against a re-simulated node to a relative error of 1.8e-16}
\expandafter\def\csname fignote@B1\endcsname{V$_t$ at I$_{cc}$ = 1e-08 A/sheet: +0.363 V erased, -0.932 V programmed one free width scalar (0.0065) fitted on the erased saturation above 2.5 V}
\expandafter\def\csname fignote@B3\endcsname{V$_t$ at 100 nA/$\mu$m; I$_{on}$ at 0.4 V overdrive}
\expandafter\def\csname fignote@B5\endcsname{erased: RMS 0.478 dec, median -0.026, n = 12 programmed: RMS 0.295 dec, median -0.011, n = 8 shaded: $\pm$0.3 decade}
\expandafter\def\csname fignote@C1\endcsname{saturated loop ($\pm$5.00 V): P$_s$ = 39.99, |P$_r$| = 31.99 $\mu$C/cm$^2$, F$_c$ = 1.400 MV/cm}
\expandafter\def\csname fignote@C2\endcsname{The $\pm$4 V gate sweep reaches 0.773 MV/cm, 0.55 of the material coercive field F$_c$ = 1.400 MV/cm measured in C1.}
\expandafter\def\csname fignote@D1\endcsname{V$_t$ at I$_{cc}$ = 0.01 $\mu$A/$\mu$m SS = 85.5 / 63.5 mV/dec (erased / programmed), steepest fit spanning $\geq$2 decades}
\expandafter\def\csname fignote@D2\endcsname{read at V$_G$ = 0.2 V linear-region correlation r = 0.940$-$1.000 over V$_{DS}$ $\leq$ 0.1 V}
\expandafter\def\csname fignote@E1\endcsname{$\Delta$P over the train = -1.858 $\mu$C/cm$^2$; shaded bands are the V$_G$ = 2 V write windows. Median gate conduction between pulses: 26.7 pA.}
\expandafter\def\csname fignote@E2\endcsname{Switching per pulse peaks at pulse 8, then falls to 0.77 of the peak by pulse 15.}
\expandafter\def\csname fignote@E6\endcsname{V$_{pgm}$ = 2 V, 15 pulses peak 11.0 aJ at pulse 8, last pulse 8.9 aJ whole train 123.1 aJ = 0.1231 fJ}
\expandafter\def\csname fignote@F1\endcsname{15 pulses, 3.79 decades 8 separable open-loop (3.0 bit) $\approx$16 with write-verify (4.0 bit)}
\expandafter\def\csname fignote@F2\endcsname{A$_{LTP}$ = -0.516, R$^2$ = 0.994 10--90 \% span = 3.6 pulses max residual 0.059, 12 of 15 points below the fit}
\expandafter\def\csname fignote@F5\endcsname{window per cycle: 11,650$\times$, 4,185$\times$, 3,577$\times$ ceiling moves 0.93$\times$, floor moves 3.0$\times$}
\expandafter\def\csname fignote@F9\endcsname{At the top level 350 K is 5.9$\times$ below 300 K while 250 K is 1.10$\times$ it -- the ladder is not monotonic in temperature, so no activation energy is extracted.}
\expandafter\def\csname fignote@G1\endcsname{V$_{pgm}$ = 2 V, 100 ns pulses rest current 0.00821 $\mu$A/$\mu$m (dotted) threshold 0.1 $\mu$A/$\mu$m (dashed), a design choice contrast at fire: 18.4$\times$ rest}
\expandafter\def\csname fignote@G2\endcsname{Dashed line: the 0.05 $\mu$C/cm$^2$ state threshold. An open marker with an arrow means the threshold was not reached within the pulses that were run.}
\expandafter\def\csname fignote@G3\endcsname{Crosses threshold on pulse 9 at V$_{pgm}$ = 2 V; the dashed line is the 0.1 $\mu$A/$\mu$m read threshold.}
\expandafter\def\csname fignote@H1\endcsname{settling: t < 32 $\mu$s = 3 $\tau_P$ plateau 4.66 $\mu$A/$\mu$m $\pm$1 \% band drift on the plateau -0.82 \%}
\expandafter\def\csname fignote@H4\endcsname{300 K, programmed: settled at 50 $\mu$s = 5 $\tau_P$, flat to 0.62 \% over 2.3 decades 400 K, programmed: settled at 37 $\mu$s = 4 $\tau_P$, flat to 0.96 \% over 2.4 decades 300 K, LTP level 8: settled at 55 $\mu$s = 6 $\tau_P$, flat to 0.62 \% over 1.3 decades}
\expandafter\def\csname fignote@H8\endcsname{measured from 1e+04 reads onward +0.0000 \% over 9.9e+05 reads apparent loss if read unsettled: -89.1 \% ($\tau_P$ transient, not disturb)}
\expandafter\def\csname fignote@I1\endcsname{The window improves 37,015$\times$ over 6 turns; turns 5 and 6 change nothing.}
\expandafter\def\csname fignote@I2\endcsname{1.0 nm buys 76$\times$ more window than 1.5 nm}
\expandafter\def\csname fignote@I3\endcsname{marker area $\propto$ gate charge 7 nm: 2.0$\times$ the window of the best higher-voltage point, at 2.5 V and the smallest gate charge (0.86 of the largest)}
\expandafter\def\csname fignote@I5\endcsname{}
\expandafter\def\csname fignote@J1\endcsname{Each axis is normalized to the better of the two devices, so outward is better. The level axis uses one shared readout-separation rule ($\geq$1.5$\times$) so the two devices are comparable; this device's headline open-loop count under its own measured cycle-to-cycle noise is smaller (see F1).}
\expandafter\def\csname fignote@J5\endcsname{1 of 15 surveyed devices report both quantities; 9 report a level count only and are drawn as left-pointing markers at an arbitrary energy. Open markers are points whose source number could not be confirmed in the paper text.}
